\section{Yu--Wang--Samworth as a statistical specialization}
\label{sec:yws}

YWS \citep{yu2015davis} gives a finite-dimensional population-gap version of
subspace perturbation theory.  Let
$\sym{Sigma},\sym{SigmaHat}\in\mathbb{R}^{p\times p}$ be symmetric, with population
eigenvalues $\lambda_1\ge\cdots\ge\lambda_p$.  For
$1\le r\le s\le p$, write $d=s-r+1$ and let $\sym{YWSV},\sym{YWSVhat}$ contain orthonormal
eigenvectors at indices $r,\ldots,s$.  With population gap
\[
  \sym{Delta}=
  \min\{\lambda_{r-1}-\lambda_r,\lambda_s-\lambda_{s+1}\}>0,
\]
YWS Theorem~2 gives
\begin{equation}
  \|\sin\sym{Theta}(\sym{YWSVhat},\sym{YWSV})\|_{\mathrm F}
  \le
  \frac{2\min\{\sqrt d\,\|\sym{SigmaHat}-\sym{Sigma}\|_{\mathrm{op}},
                    \|\sym{SigmaHat}-\sym{Sigma}\|_{\mathrm F}\}}
       {\sym{Delta}}.
  \label{eq:yws-sintheta}
\end{equation}
An aligned-frame bound follows with an additional factor $\sqrt2$.
Here $\lambda_0=+\infty$ and $\lambda_{p+1}=-\infty$ are the usual
boundary conventions. Theorem~3 applies the same population-gap principle to left and right singular
subspaces through the corresponding Gram matrices.  The formalization checks
Theorems~1--3 and the supporting identities and conventions needed for those
claims; Appendix~\ref{app:yws-coverage} gives the exact scope.

The comparison also requires attention to the formulas and boundary
conventions in the archived 2015 article. Appendix~\ref{app:yws-coverage}
gives the checked corrections and distinguishes published source wording
from the mathematical hypotheses of the Lean theorems.
